\section{From posterior estimates to a spectral gap}\label{sec:observation}

We now prove the transfer announced in the introduction. Its inputs are
a covariance bound, a square-reweighted mean bound, and a terminal gap.
The stopping rule and good-disorder event must be independent of the test
function, so we keep the exceptional-path estimate explicit. The argument
uses the Gaussian-channel formulation of localization
\cite{ElAlaouiMontanari2021} and the variance/Dirichlet-form framework
of~\cite{CE2025}; the finite-spin identities are derived here.

\subsection{Posterior martingales and Dirichlet forms}
Let $\mathcal F_t=\sigma(Y(s):s\le t)$ and write $\E_t$ for expectation under
$\mu_{Y(t)}$. Bayes' formula gives this posterior: the likelihood of an
observation path given $X=x$ is proportional to
$\exp(x^{\mathsf T}Y(t)-tn/2)$. The process
\[
 I(t)=Y(t)-\int_0^t\E_s x\,\dd s
\]
is the innovation Brownian motion in $\mathcal F_t$: conditional
expectation makes it a continuous martingale, and its quadratic
covariation is $tI_n$. The Brownian martingale characterization therefore
applies. Applying It\^o's formula to the finitely many posterior weights
gives, for every spin function $g$,
\begin{equation}\label{eq:filtering}
 \dd(\E_t g)=\Cov_t(g,x)^{\mathsf T}\dd I(t).
\end{equation}
To verify the cancellation, differentiate the posterior expectation with
respect to $Y$; its gradient is $\Cov_t(g,x)$. Each likelihood weight
$\exp(x^{\mathsf T}Y(t)-tn/2)$ satisfies a stochastic exponential
equation. The quotient rule for their weighted sums then gives
\eqref{eq:filtering}: subtracting $\E_t x\,\dd t$ from $\dd Y$ removes
the finite-variation term.

At every deterministic time,
\begin{equation}\label{eq:dirichlet-monotonicity}
 \E\D_{Y(t)}(f)\le\D_0(f).
\end{equation}
This is conditional-variance contraction: the $i$th summand on the left
is $\E\Var(f(X)\mid X_{-i},\mathcal F_t)$, and the one on the right is
$\E\Var(f(X)\mid X_{-i})$. The conditional variance decomposition gives
the inequality after averaging over observations. In this section, an
expectation without a field subscript refers to that experiment, with
$J$ held fixed.

\subsection{Stopping before stability fails}
Stop before the first possible loss of the field estimates. More precisely,
fix the terminal horizon $T$, take the observation stopping time $\tau_0$
from Corollary~\ref{cor:stopping}, and set $\tau=T\wedge\tau_0$. This
rule does not use the test function. Before $\tau$, stability holds through
residual threshold $\rho_0\sqrt n$, so both posterior estimates apply.
On the specified matrix event, $\Prob(\tau_0\le T)\le e^{-an}$.

For a fixed $J$, choose the canonical terminal good set
\[
 \mathcal T_n(J)=\{h:\Var_{\mu_h}(g)\le2\D_h(g)\text{ for every }g\}.
\]
Membership means positivity of a finite-dimensional symmetric quadratic
form with coefficients continuous in $h$, so this set is Borel.
Proposition~\ref{prop:terminal-gap} controls the probability that $Y(T)$
misses it. Declare the stopped path bad if $\tau_0\le T$, or if
$\tau_0>T$ and $Y(T)\notin\mathcal T_n(J)$. The event is
$\mathcal F_\tau$-measurable: before $T$ only stopping can make it bad,
and at $T$ the remaining condition is determined by $Y(T)$. The two
failure probabilities give, on the typical matrix events,
\begin{equation}\label{eq:ordinary-bad}
 \Prob(\mathrm{bad})\le 2e^{-a n}
\end{equation}
for a fixed $a>0$, after decreasing the rate if necessary.

Suppose $\E_{\mu_0}f=0$ and $\E_{\mu_0}f^2=1$. Set
\[
 M_t=\E_t f,\quad N_t=\E_t f^2,\quad V_t=N_t-M_t^2,
 \quad d_0=\D_0(f).
\]
By \eqref{eq:filtering} and optional stopping, $u(t)=\E V_{t\wedge\tau}$ is
absolutely continuous and satisfies, almost everywhere,
\[
 u'(t)=-\E\bigl[\one_{\{t<\tau\}}\norm{\Cov_t(f,x)}^2\bigr].
\]
The directional estimate bounds the loss term before stopping. Together
with \eqref{eq:dirichlet-monotonicity}, applied at the deterministic time
$t$, it gives $u'(t)\ge-C_{\mathrm d}(u(t)+d_0)$, with fixed
$C_{\mathrm d}$. Solving this scalar inequality from $u(0)=1$ yields
\begin{equation}\label{eq:variance-lower}
 \E V_\tau\ge v-(1-v)d_0,\qquad v=e^{-C_{\mathrm d}T}>0.
\end{equation}
All uses of optional stopping here are at bounded times and for bounded
random variables at each fixed $n$.

\subsection{Changing the law of the observations}
The variance lower bound alone does not control how much variance lies on
bad paths. For that purpose change the initial spin law to $f^2\mu_0$,
which is a probability law by normalization, and call the resulting
observation law $\Prob^f$. Its density on $\mathcal F_t$ is $N_t$, and
its spin posterior is $\nu_t=f^2\mu_{Y(t)}/N_t$. Since $\mu_{Y(t)}$ has
full support and $f$ is nonzero, $N_t>0$. Bounded stopping gives the
density $N_\tau$ on $\mathcal F_\tau$.
Let
\[
 a_t=\E_{\nu_t}x-\E_t x.
\]
Equation~\eqref{eq:filtering} with $g=f^2$ gives
$\dd N_t/N_t=a_t^{\mathsf T}\dd I(t)$. Under $\Prob^f$, $I$ has drift $a_t$.
The drift-energy entropy identity, in the form discussed
by Lehec~\cite{Lehec2013}, has a direct finite-mixture proof here.
It\^o's formula for $\log N_t$ yields
\begin{equation}\label{eq:stopped-entropy}
 D\bigl(\Prob^f|_{\mathcal F_\tau}\,\Vert\,
        \Prob|_{\mathcal F_\tau}\bigr)
 =\E^f\log N_\tau
 =\frac12\E^f\int_0^\tau\norm{a_t}^2\,\dd t.
\end{equation}
The stopped identity needs no limiting integrability argument. At fixed
$n$, the drift satisfies $\norm{a_t}\le2\sqrt n$ and $\tau\le T$.
Novikov's condition therefore holds, and the stochastic integral in the
log-density formula has zero mean.

Before stopping, applying \eqref{eq:input-weighted} to $f$ and to $1$ gives
\[
 b_t:=\frac{\norm{a_t}}{\sqrt n}
 \le2C_{\mathrm w}\delta+\frac{C_\delta}{\sqrt n}
                   \left(2+\frac{\D_{Y(t)}(f)}{N_t}\right).
\]
Since $b_t\le2$, we may use $b_t^2\le2b_t$. At deterministic $t$, the density
identity and \eqref{eq:dirichlet-monotonicity} imply
\begin{equation}\label{eq:weighted-dirichlet}
 \E^f\frac{\D_{Y(t)}(f)}{N_t}
 =\E\D_{Y(t)}(f)\le d_0.
\end{equation}
Write $\mathcal H_\tau$ for the entropy in \eqref{eq:stopped-entropy}.
The preceding estimates give the quantitative bound
\begin{equation}\label{eq:entropy-bound}
 \frac{\mathcal H_\tau}{n}
 \le 2C_{\mathrm w}T\delta
       +\frac{C_\delta T}{\sqrt n}(2+d_0).
\end{equation}
The separation of constants is essential: $C_{\mathrm w}$ is independent
of $\delta$, whereas $C_\delta$ may grow when $\delta$ decreases.
We next convert this entropy bound into a bound on the bad-path probability.
If an event has probabilities $p$ and $q$ under two laws, the log-sum
inequality for the event and its complement bounds relative entropy below
by $p\log(1/q)-\log2$. With $p=\Prob^f(\mathrm{bad})$,
\eqref{eq:ordinary-bad} therefore gives, for large $n$,
\[
 p\le\frac{\mathcal H_\tau+\log2}{an-\log2}.
\]
If the ordinary bad probability is zero, absolute continuity gives $p=0$.
For $d_0\le1$, \eqref{eq:entropy-bound} therefore implies
\begin{equation}\label{eq:weighted-bad-bound}
 p\le\frac{4C_{\mathrm w}T}{a}\delta
       +\frac{6C_\delta T}{a\sqrt n}
       +\frac{2\log2}{an}.
\end{equation}
Choose $\delta>0$ once so that the first term is at most $v/8$.
With this choice fixed, enlarge $n$ until the last two terms sum to at
most $v/8$. The resulting bound holds simultaneously for all normalized
$f$ with $d_0\le1$:
\begin{equation}\label{eq:weighted-bad}
 \Prob^f(\mathrm{bad})\le v/4.
\end{equation}

Figure~\ref{fig:observation-flow} records how the two posterior estimates
will control the stopped variance and its exceptional contribution.
\begin{figure}[htbp]
\centering
\small
\setlength{\tabcolsep}{5pt}
\renewcommand{\arraystretch}{1.25}
\begin{tabular}{c@{\quad}c}
\fbox{\begin{minipage}[c]{0.36\linewidth}\centering
Ordinary rarity\\[3pt]
$\Prob(\mathrm{bad})\le2e^{-an}$
\end{minipage}}
&\fbox{\begin{minipage}[c]{0.51\linewidth}\centering
Weighted-mean estimate\\[3pt]
$\displaystyle\frac{\mathcal H_\tau}{n}
 \le2C_{\mathrm w}T\delta
       +\frac{C_\delta T}{\sqrt n}(2+d_0)$
\end{minipage}}\\[4pt]
\multicolumn{2}{c}{\shortstack{\(\Downarrow\)\quad entropy inequality;
 choose $\delta$, then $n$}}\\[3pt]
\multicolumn{2}{c}{\fbox{\begin{minipage}{0.89\linewidth}\centering
Variance on bad paths\\[3pt]
$\displaystyle
 \E[\one_{\mathrm{bad}}V_\tau]
 \le\E[\one_{\mathrm{bad}}N_\tau]
 =\Prob^f(\mathrm{bad})\le v/4$
\end{minipage}}}\\[9pt]
&\shortstack{\(\Downarrow\)\quad terminal gap on good paths\\[2pt]
 $\E[\one_{\mathrm{good}}V_\tau]\le2d_0$}\\[4pt]
\fbox{\begin{minipage}[c]{0.36\linewidth}\centering
Directional covariance estimate\\[3pt]
$\E V_\tau\ge v-(1-v)d_0$
\end{minipage}}
&\fbox{\begin{minipage}[c]{0.51\linewidth}\centering
Total remaining variance\\[3pt]
$\E V_\tau\le2d_0+v/4$
\end{minipage}}\\[4pt]
\multicolumn{2}{c}{$\displaystyle
 v-(1-v)d_0\le2d_0+v/4\qquad\Longrightarrow\qquad d_0\ge v/4$}
\end{tabular}
\caption{Variance retention and control of exceptional paths, for a fixed
typical matrix. Normalize $f$ by $\E_{\mu_0}f=0$ and
$\E_{\mu_0}f^2=1$, and put $d_0=\D_0(f)\le1$.
Along $Y(t)=tX+B(t)$, let $V_t=\Var_{\mu_{Y(t)}}f$,
$N_t=\E_{\mu_{Y(t)}}f^2$, and $\tau=T\wedge\tau_0$.
Here $\mathcal H_\tau$ is the relative entropy of the stopped observation
law induced by the spin law $f^2\mu_0$, relative to the ordinary law,
and $v=e^{-C_{\mathrm d}T}$.
The density $N_\tau$ converts its bad-path probability into an upper
bound on exceptional variance. If $d_0>1$, the final lower bound is automatic.}
\label{fig:observation-flow}
\end{figure}


\subsection{Completion of the theorem}
We can now close the argument, keeping the order of choices explicit.
First choose $T$ from the terminal estimate. Next choose the stable-field
constants, the ordinary failure rate $a$, and the finite diagnostic length
for the directional estimate; these choices determine $C_{\mathrm d}$
and $v$. Root stability and the operator bound on $J$ determine the
leading weighted-mean constant $C_{\mathrm w}$ independently of selection
depth. We may therefore select $\delta$ in
\eqref{eq:weighted-bad-bound}, then its diagnostic length and
$C_\delta$, and finally $n_0$. Every depth and tolerance is fixed before
$n$ grows.

With the choices made, the proof needs a finite intersection of matrix
events. Include the directional word event and the word events for the
selected weighted-mean construction, each with its fixed length and
coefficient bounds. Lemma~\ref{lem:words} accommodates each fixed
$M_0$, and Lemma~\ref{lem:implicit-construction} supplies an actual
inverse margin independent of depth. Add the events of
Propositions~\ref{prop:stable-fields} and~\ref{prop:terminal-gap}.
Their intersection still has probability tending to one, and all constants
and the threshold $n_0$ are uniform on it and over the test functions.

For such a matrix and a normalized $f$ with $d_0\le1$, use $V_\tau\le N_\tau$
and the stopped density identity to obtain
\[
 \E[\one_{\mathrm{bad}}V_\tau]
 \le\E[\one_{\mathrm{bad}}N_\tau]
 =\Prob^f(\mathrm{bad})\le v/4.
\]
On a good path $\tau=T$ and the terminal gap applies, so
\[
 \E[\one_{\mathrm{good}}V_\tau]
 \le2\E\D_{Y(T)}(f)\le2d_0.
\]
Combining these estimates with \eqref{eq:variance-lower} gives
\[
 v-(1-v)d_0\le2d_0+v/4,
 \qquad\text{hence}\qquad
 d_0\ge\frac{3v}{4(3-v)}\ge\frac v4.
\]
The case $d_0>1$ already satisfies this bound. Center and rescale any
nonconstant function to obtain Theorem~\ref{thm:main} once the three
inputs have been proved, with the finite constant
$C_\beta=4e^{C_{\mathrm d}T}$. We do not optimize this constant.
For the clock conversion, the self-adjoint operator
$H=\sum_i(I-P_i^0)$ has form $\D_0$, and one discrete attempt is
$I-H/n$. Its gap is therefore the asserted unscaled gap divided by $n$.
